\documentclass[11pt,a4paper]{article}
\usepackage[utf8]{inputenc}
\usepackage[T1]{fontenc}
\usepackage{lmodern}
\usepackage{textcomp}
\usepackage{amsmath}
\usepackage{array}
\usepackage[hmargin=1.6cm,vmargin=1.2cm]{geometry}
\usepackage{enumitem}
\pagestyle{empty}
\setlength{\parindent}{0pt}
\setlength{\parskip}{0pt}
\begin{document}
\raggedright
{\Large\bfseries Tarea de pensamiento computacional}
\vspace{10pt}
\begin{enumerate}[label=\textbf{\arabic*.},leftmargin=*,itemsep=6pt,topsep=0pt]
\item \begin{minipage}[t]{\linewidth}
¿Qué frase resume mejor el pensamiento computacional según el aula?

\vspace{4pt}
\begin{tabular}{@{}>{\raggedright\arraybackslash}p{0.47\linewidth}>{\raggedright\arraybackslash}p{0.47\linewidth}@{}}
a)~Un proceso mental para formular problemas y diseñar soluciones que pueda ejecutar un agente que procesa información & b)~La habilidad de escribir código en un lenguaje de programación \\
c)~El uso de computadoras para buscar información en internet & d)~Una técnica exclusiva de los ingenieros en sistemas \\
\end{tabular}
\end{minipage}

\item \begin{minipage}[t]{\linewidth}
Una frase dice que el pensamiento computacional implica resolver problemas, diseñar sistemas y comprender el comportamiento humano. ¿De quién es?

\vspace{4pt}
\begin{tabular}{@{}>{\raggedright\arraybackslash}p{0.47\linewidth}>{\raggedright\arraybackslash}p{0.47\linewidth}@{}}
a)~Manuel Castells, 2001 & b)~Seymour Papert, 2006 \\
c)~Jeannette Wing, 2006 & d)~Jeannette Wing, 2011 \\
\end{tabular}
\end{minipage}

\item \begin{minipage}[t]{\linewidth}
Para estudiar 10 temas, María los separa por dificultad y los reparte entre los días que faltan para el examen. ¿Qué pilar aplica?

\vspace{4pt}
\begin{tabular}{@{}>{\raggedright\arraybackslash}p{0.47\linewidth}>{\raggedright\arraybackslash}p{0.47\linewidth}@{}}
a)~Pensamiento algorítmico & b)~Abstracción \\
c)~Descomposición & d)~Reconocimiento de patrones \\
\end{tabular}
\end{minipage}

\item \begin{minipage}[t]{\linewidth}
En la calculadora, el mismo esquema número, operador, número e igual sirve para sumar, restar, multiplicar o dividir. ¿Qué pilar se ve?

\vspace{4pt}
\begin{tabular}{@{}>{\raggedright\arraybackslash}p{0.47\linewidth}>{\raggedright\arraybackslash}p{0.47\linewidth}@{}}
a)~La salida del modelo Entrada, Proceso y Salida & b)~Descomposición \\
c)~Abstracción & d)~Reconocimiento de patrones \\
\end{tabular}
\end{minipage}

\item \begin{minipage}[t]{\linewidth}
En una aplicación de recetas, cada receta se reduce a una lista de ingredientes y un tiempo, y lo demás se descarta. ¿Qué pilar es?

\vspace{4pt}
\begin{tabular}{@{}>{\raggedright\arraybackslash}p{0.47\linewidth}>{\raggedright\arraybackslash}p{0.47\linewidth}@{}}
a)~Descomposición & b)~Abstracción \\
c)~Reconocimiento de patrones & d)~Pensamiento algorítmico \\
\end{tabular}
\end{minipage}

\item \begin{minipage}[t]{\linewidth}
Un algoritmo recibe 4 y 5 y entrega 9 hoy, mañana y siempre. ¿Qué característica muestra?

\vspace{4pt}
\begin{tabular}{@{}>{\raggedright\arraybackslash}p{0.47\linewidth}>{\raggedright\arraybackslash}p{0.47\linewidth}@{}}
a)~Determinístico, porque da el mismo resultado con los mismos datos & b)~Efectivo, porque no llega a ninguna solución \\
c)~Finito, porque nunca termina & d)~Preciso, porque cambia de resultado cada día \\
\end{tabular}
\end{minipage}

\item \begin{minipage}[t]{\linewidth}
En el flujograma para revisar el wifi, la pregunta «¿Fallan todos los dispositivos?» se dibuja con

\vspace{4pt}
\begin{tabular}{@{}>{\raggedright\arraybackslash}p{0.47\linewidth}>{\raggedright\arraybackslash}p{0.47\linewidth}@{}}
a)~un óvalo, porque es el inicio o el fin & b)~un rectángulo, porque es un proceso \\
c)~un rombo, porque es una decisión de sí o no & d)~una flecha, porque indica el flujo \\
\end{tabular}
\end{minipage}

\item \begin{minipage}[t]{\linewidth}
En un cajero automático, ¿cuál es una entrada?

\vspace{4pt}
\begin{tabular}{@{}>{\raggedright\arraybackslash}p{0.47\linewidth}>{\raggedright\arraybackslash}p{0.47\linewidth}@{}}
a)~La comparación del monto con el saldo & b)~Los billetes entregados \\
c)~El recibo impreso & d)~La contraseña que escribe el usuario \\
\end{tabular}
\end{minipage}

\item \begin{minipage}[t]{\linewidth}
En una calculadora escribes 8 $\times$ 7 y pulsas igual. ¿Qué es la salida?

\vspace{4pt}
\begin{tabular}{@{}>{\raggedright\arraybackslash}p{0.47\linewidth}>{\raggedright\arraybackslash}p{0.47\linewidth}@{}}
a)~56 en la pantalla & b)~15 en la pantalla \\
c)~El cálculo de la multiplicación & d)~Los números 8 y 7 \\
\end{tabular}
\end{minipage}

\item \begin{minipage}[t]{\linewidth}
Juan compra repuestos para su bicicleta sin revisar cuál es la falla. Los pone y la bicicleta sigue sin funcionar. ¿Qué error cometió según la metodología?

\vspace{4pt}
\begin{tabular}{@{}>{\raggedright\arraybackslash}p{0.47\linewidth}>{\raggedright\arraybackslash}p{0.47\linewidth}@{}}
a)~Dedicó demasiado tiempo a verificar & b)~Saltó a ejecutar sin comprender el problema \\
c)~Descompuso el problema en partes & d)~Comprendió el problema antes de diseñar \\
\end{tabular}
\end{minipage}

\end{enumerate}
\end{document}
